\documentclass[aspectratio=169,11pt]{beamer}

\usetheme{Madrid}
\usecolortheme{default}
\usefonttheme{professionalfonts}
\setbeamertemplate{navigation symbols}{}
\setbeamertemplate{footline}[frame number]

\usepackage{amsmath, amssymb, booktabs}

\title{Social Norms, Bargaining, And Institutions Shaping Gendered Work}
\subtitle{Gender Study --- Week 5}
\author{Teaching deck draft}
\date{}

\begin{document}

\begin{frame}
  \titlepage
\end{frame}

\begin{frame}{Central question}
\begin{itemize}
    \item How do norms, bargaining power, and institutions shape gendered work even when prices and formal rules change?
    \item Norms matter for labor when they change participation, hours, specialization, search, commuting, competition, promotion, retention, or job quality.
    \item Week 5 discipline: treat norms as identifiable mechanisms, not residual explanations.
\end{itemize}
\end{frame}

\begin{frame}{Why norms and institutions are labor objects}
\begin{itemize}
    \item They shift feasible jobs, acceptable hours, mobility, bargaining positions, resource control, and authority.
    \item They alter responses to wages, childcare, leave, safety, transparency, and legal rights.
    \item The same average gap can come from prices, law, selection, norms, bargaining, or their interaction.
\end{itemize}
\end{frame}

\begin{frame}{Week 2 $\rightarrow$ Week 5 bridge}
\small
\begin{tabular}{p{0.28\linewidth} p{0.58\linewidth}}
\toprule
Week 2 object & Week 5 mechanism \\
\midrule
Household allocation & bargaining, control over resources, threat points \\
Care work and child penalties & caregiver norms and acceptable-hours expectations \\
Fertility timing & work-family beliefs and marriage-market expectations \\
Labor supply & wage response filtered through sanctions and household rules \\
\bottomrule
\end{tabular}
\end{frame}

\begin{frame}{Bridge from Weeks 3 and 4}
\begin{itemize}
    \item Week 3: aspirations, field choice, and early careers can respond to anticipated marriage-market and work-family norms.
    \item Week 4: firms reward self-promotion, competition, social access, long hours, mobility, and negotiation.
    \item Week 5 connects those worker and firm margins to informal institutions and bargaining power.
\end{itemize}
\end{frame}

\begin{frame}{Formal versus informal institutions}
\small
\begin{tabular}{p{0.28\linewidth} p{0.58\linewidth}}
\toprule
Institution type & Examples \\
\midrule
Formal & law, property rights, contract rules, leave, tax, childcare, pay rules, legal protection \\
Informal & breadwinning, caregiving, acceptable jobs, mobility, competition, self-promotion, authority \\
\bottomrule
\end{tabular}
\vspace{0.5em}
\begin{itemize}
    \item Formal rules can codify norms or relax them.
    \item Informal expectations can make formal rights more or less usable.
\end{itemize}
\end{frame}

\begin{frame}{Five claims to keep separate}
\begin{enumerate}
    \item Price claim: wages, childcare, taxes, or commute costs moved.
    \item Legal claim: a formal rule permits, prohibits, or protects behavior.
    \item Selection claim: people with different preferences or constraints sort.
    \item Norm claim: expectations, sanctions, identity, or coordination change payoffs.
    \item Bargaining claim: threat points or control over resources change allocation.
\end{enumerate}
\end{frame}

\begin{frame}{Methodological taxonomy of norms}
\small
\begin{tabular}{p{0.25\linewidth} p{0.46\linewidth} p{0.18\linewidth}}
\toprule
Norm channel & Mechanism & Margin \\
\midrule
Identity / preference & utility cost of violating role & hours, specialization \\
Sanction belief & perceived approval, reputation, status & search, work entry \\
Coordination & stable social equilibrium & education, marriage, jobs \\
Bargaining & threat point and resource control & labor supply, savings \\
Embedded in rules & formal institution carries norm & rights, mobility, pay \\
\bottomrule
\end{tabular}
\end{frame}

\begin{frame}{A household surplus object}
\[
S(h_w,h_m,q)
=
U_w(c,\ell_w,q)
+
U_m(c,\ell_m,q)
-
\kappa N(h_w,h_m,q;g)
-
C(h_w,h_m,q)
\]
\begin{itemize}
    \item \(N(\cdot;g)\): norm-violation cost under gendered rule \(g\).
    \item \(\kappa\): salience or expected sanction cost.
    \item A higher female wage can have a muted response if the norm penalty is high.
\end{itemize}
\end{frame}

\begin{frame}{Bargaining with a gender margin}
\[
\max_{h_w,h_m,q}
\left[V_w(h_w,h_m,q)-T_w(g,F)\right]^{\theta_w}
\left[V_m(h_w,h_m,q)-T_m(g,F)\right]^{\theta_m}
\]
\begin{itemize}
    \item \(T_w(g,F)\): woman's threat point under norms \(g\) and formal institutions \(F\).
    \item Norms can lower threat points through stigma, mobility limits, or weak control over earnings.
    \item Formal institutions can raise threat points through rights, accounts, protection, and transfers.
\end{itemize}
\end{frame}

\begin{frame}{Breadwinner, caregiving, and marriage-market norms}
\begin{itemize}
    \item Breadwinner norms can affect relative income, labor supply, and home production.
    \item Caregiving norms can make long hours or uninterrupted careers socially costly.
    \item Marriage-market norms can alter education, field choice, ambition, and early career investments.
    \item The observed margin must be named: earnings share, hours, reported ambition, or investment.
\end{itemize}
\end{frame}

\begin{frame}{Misperceptions and local sanctions}
\begin{itemize}
    \item Workers and spouses may privately support women's work but misperceive community approval.
    \item Information treatments identify belief-about-acceptability channels when job-search behavior moves.
    \item Status and reputation concerns can change labor supply even without a wage or legal shock.
    \item Key diagnostic: does the design shift beliefs, sanctions, or prices?
\end{itemize}
\end{frame}

\begin{frame}{Competition, authority, and self-promotion}
\begin{itemize}
    \item Competitive job design can change who applies before the firm evaluates anyone.
    \item Authority and leadership norms shape self-promotion, negotiation, visibility, and willingness to compete.
    \item Overwork norms interact with caregiving norms in professional careers.
    \item Entry, promotion, and retention are distinct margins.
\end{itemize}
\end{frame}

\begin{frame}{Acceptable jobs, commuting, and mobility}
\begin{itemize}
    \item Mobility constraints can reflect safety, household bargaining, local reputation, or acceptable-job norms.
    \item Revealed-preference search data can estimate wage-commute tradeoffs.
    \item A commute gap is not automatically a taste parameter.
    \item Ask what information is observed: offers, applications, accepted jobs, commutes, wages, or household constraints.
\end{itemize}
\end{frame}

\begin{frame}{Formal laws interacting with norms}
\begin{itemize}
    \item Gendered laws can restrict employment, mobility, property, entrepreneurship, or workplace rights.
    \item Formal equality can be weak if enforcement is not credible or reporting is costly.
    \item Legal rights can also change informal expectations by changing feasible behavior and threat points.
    \item Cross-country legal evidence must be read with enforcement and norm interaction in mind.
\end{itemize}
\end{frame}

\begin{frame}{Empirical designs and what they identify}
\scriptsize
\setlength{\tabcolsep}{3pt}
\begin{tabular}{p{0.24\linewidth} p{0.28\linewidth} p{0.30\linewidth}}
\toprule
Design & Identifying variation & Best observed margin \\
\midrule
Second-generation / ancestry & origin-linked norms in common destination & work, fertility, beliefs \\
Historical origins & long-run production or institutions & persistence of gender roles \\
Norm intervention & randomized information about support & job search, take-up, beliefs \\
Household bargaining & resource control or threat-point shift & labor supply, allocation \\
Job-entry experiment & compensation or public/private elicitation & applications, ambition \\
Commuting/search & revealed job choices & commute-wage tradeoff \\
Legal institutions & law variation over time/countries & participation, rights, enforcement \\
\bottomrule
\end{tabular}
\end{frame}

\begin{frame}{What each anchor paper teaches}
\scriptsize
\begin{tabular}{p{0.31\linewidth} p{0.24\linewidth} p{0.29\linewidth}}
\toprule
Anchor & Labor margin & Mechanism lesson \\
\midrule
Fernandez--Fogli; Alesina--Giuliano--Nunn & work, fertility, roles & culture and persistence \\
Bursztyn--Gonzalez--Yanagizawa-Drott & job matching & misperceived acceptability \\
Bertrand--Kamenica--Pan & earnings, hours, home production & breadwinner identity \\
Bursztyn--Fujiwara--Pallais & ambitions, investment & marriage-market incentives \\
Flory--Leibbrandt--List & applications & competitive job design \\
Le Barbanchon--Rathelot--Roulet & search, commuting & mobility tradeoffs \\
Hyland--Djankov--Goldberg & workforce outcomes & formal legal institutions \\
\bottomrule
\end{tabular}
\end{frame}

\begin{frame}{Week 5 lab}
\begin{itemize}
    \item Reproduce: synthetic norm-misperception information exercise inspired by Bursztyn, Gonzalez, and Yanagizawa-Drott.
    \item Diagnose: classify outputs as beliefs, sanctions, bargaining, search, or selection.
    \item Transfer: synthetic relative-income threshold exercise inspired by Bertrand, Kamenica, and Pan.
    \item Optional memo: financial control or commuting revealed-preference designs.
\end{itemize}
\end{frame}

\begin{frame}{Bridge to Week 6 policy incidence}
\begin{itemize}
    \item Week 5: norms are rules of conduct, bargaining environments, and informal institutions.
    \item Week 6: policies change prices and formal rules, then interact with those norms.
    \item Childcare, leave, taxation, quotas, transparency, safety, and protection can all have norm-mediated incidence.
\end{itemize}
\end{frame}

\begin{frame}{Takeaways}
\begin{itemize}
    \item Norms are not a catch-all residual; they are mechanisms with empirical signatures.
    \item Separate norms from prices, law, selection, bargaining, and firm design.
    \item Formal and informal institutions interact through feasibility, enforcement, information, and threat points.
    \item Strong evidence names the variation, the margin, and the specific channel.
\end{itemize}
\end{frame}

\end{document}
